\chapter{Моделювання загроз, ланцюги ураження та оцінювання ризиків злоякісного коду}
\chaptershort{Моделювання загроз і оцінювання ризиків}
\label{chap:ch04}

Нормативні вимоги (розділ~\ref{chap:ch02}) визначають, що має бути захищено, а теоретичні основи (розділ~\ref{chap:ch03}) --- принципові межі виявлення. Між ними лежить процес, що пов'язує активи, загрози й заходи в систему обґрунтованих рішень: моделювання загроз і управління ризиками. Для протидії злоякісному коду цей процес має специфіку: загроза реалізується як багатоетапна операція (ланцюг ураження), а не одномоментна подія; шкода залежить від типу коду (крадіжка проти шифрування проти знищення); частина атак використовує самопоширення; значна частина ризику пов'язана з ланцюгом постачання й людським чинником. Розділ розглядає процес управління ризиками, методології моделювання загроз (ланцюги ураження, бази ATT\&CK, дерева атак), кількісне оцінювання ризику й моделювання поширення, пріоритизацію заходів та економічне обґрунтування, і завершується кейсом побудови моделі загроз організації. Матеріал пов'язує технічні розділи (\ref{chap:ch05}--\ref{chap:ch15}) з операційними (\ref{chap:ch16}--\ref{chap:ch17}).

\section{Процес управління ризиками}\label{sec:ch04-process}

Управління ризиками щодо злоякісного коду здійснюється відповідно до ISO/IEC~27005~\cite{iso-27005}, ISO~31000~\cite{iso-31000} і NIST~SP~800-30~\cite{ch04:nist30}: встановлення контексту, оцінювання ризиків (ідентифікація, аналіз, оцінка), оброблення, прийняття залишкових ризиків, наскрізні комунікація й моніторинг. NIST~SP~800-30 деталізує оцінювання через ланцюг «джерело загрози $\to$ подія загрози $\to$ вразливість $\to$ вплив».

Специфіка контексту для ЗПК --- визначення релевантного \emph{профілю порушника} (розділ~\ref{chap:ch01}): модель загроз для рядової організації (масові кіберзлочинці, вимагачі) принципово відрізняється від моделі для об'єкта критичної інфраструктури чи державної установи (цільові APT, деструктивні кампанії). Активами, що потребують захисту, є не лише дані, а й доступність сервісів, цілісність систем керування, довіра до ланцюга постачання, безперервність діяльності. Критерії впливу мають охоплювати фінансові втрати, порушення функціонування, витоки персональних даних, регуляторні наслідки й, для критичної інфраструктури, вплив на безпеку людей.

Важлива особливість моделювання загроз щодо ЗПК полягає в тому, що загроза є \emph{адаптивною}: на відміну від стихійних лих чи випадкових відмов, зловмисник спостерігає за захистом і змінює тактику у відповідь. Тому статична модель загроз швидко застаріває, і процес управління ризиками має бути циклічним із регулярним переглядом на основі актуальної розвідки про загрози (threat intelligence, розділ~\ref{chap:ch17}). Другою особливістю є роль \emph{людського чинника}: більшість успішних проникнень починається з дії людини (відкриття вкладення, перехід за посиланням, встановлення ПЗ), тому модель загроз має охоплювати не лише технічні вразливості, а й соціальну інженерію, і відповідно передбачати заходи навчання й технічного обмеження наслідків людських помилок. Третьою --- \emph{каскадність} впливу: зараження одного вузла в сучасних взаємопов'язаних середовищах може швидко поширитися, тому оцінювання має враховувати не лише прямий, а й поширений вплив через сегменти мережі, довірчі зв'язки й ланцюг постачання.

\section{Моделювання загроз через ланцюг ураження}\label{sec:ch04-killchain}

\subsection{Класичний ланцюг ураження}

Модель \term{ланцюга ураження} (Cyber Kill Chain)~\cite{lockheed-kc} описує злоякісну операцію як послідовність обов'язкових етапів (рис.~\thechapter.1): розвідка, озброєння, доставка, експлуатація, встановлення, командування й керування (C2), дії за цілями. Ключова ідея інтелектуально-керованого захисту --- захиснику достатньо надійно перервати \emph{будь-який} обов'язковий етап, щоб зірвати операцію, тоді як зловмиснику потрібно успішно пройти \emph{всі}. Це створює асиметрію на користь захисника за умови глибоко ешелонованого захисту: кожен етап є можливістю для виявлення й реагування.

\begin{tikzfig}
\begin{tikzpicture}[font=\scriptsize,
  s/.style={draw,rounded corners=2pt,minimum height=6mm,text width=17mm,align=center},
  >={Stealth[length=1.5mm]}]
\node[s] (a) at (0,0) {Розвідка};
\node[s] (b) at (0,-0.85) {Озброєння};
\node[s] (c) at (0,-1.7) {Доставка};
\node[s] (d) at (0,-2.55) {Експлуатація};
\node[s] (e) at (4,0) {Встановлення};
\node[s] (f) at (4,-0.85) {C2 (керування)};
\node[s] (g) at (4,-1.7) {Дії за цілями};
\draw[->] (a)--(b); \draw[->] (b)--(c); \draw[->] (c)--(d);
\draw[->] (d.east) -- ++(0.4,0) |- (e.west);
\draw[->] (e)--(f); \draw[->] (f)--(g);
\node[s,fill=black!8] (def) at (4,-2.6) {розрив будь-якого етапу зриває операцію};
\end{tikzpicture}
\figcap{Рис.~\thechapter.1 --- Ланцюг ураження злоякісної операції (за \cite{lockheed-kc})}
\end{tikzfig}

Розвитком є \term{уніфікований ланцюг ураження}~\cite{ch04:pols}, що об'єднує класичну модель із тактиками ATT\&CK і краще відображає сучасні багатоетапні операції з рухом усередині мережі (lateral movement) між первинним проникненням і досягненням цілей. Для внутрішньомережевої фази характерний цикл: закріплення $\to$ розвідка всередині $\to$ підвищення привілеїв $\to$ рух до нових вузлів, що повторюється, поки зловмисник не досягне цілі.

Кожен етап ланцюга має власні характерні засоби протидії. На етапі \emph{доставки} діють фільтрація пошти й вебтрафіку, пісочниці для вкладень (розділ~\ref{chap:ch07}), навчання користувачів; на етапі \emph{експлуатації} --- усунення вразливостей, засоби запобігання експлойтам, принцип найменших привілеїв; на етапі \emph{встановлення} --- контроль застосунків, моніторинг персистентності (розділ~\ref{chap:ch05}), захист кінцевих точок (розділ~\ref{chap:ch10}); на етапі \emph{C2} --- мережеве виявлення, аналіз DNS і репутації, блокування відомої інфраструктури (розділ~\ref{chap:ch11}); на етапі \emph{дій за цілями} --- виявлення аномальної поведінки, обмеження доступу, захист резервних копій. Така відповідність «етап --- контрзахід» є практичним інструментом проєктування ешелонованого захисту: для кожного обов'язкового етапу має бути щонайменше один незалежний рубіж виявлення або запобігання. Концепція «воронки виявлення» розвиває цю ідею: що далі просувається зловмисник, то менше в нього залишається способів уникнути виявлення, бо кожна дія в мережі залишає сліди.

\subsection{База знань ATT\&CK як модель загроз}

База MITRE~ATT\&CK~\cite{mitre-attack,ch04:strom} систематизує реальну поведінку зловмисників за тактиками (цілі: початковий доступ, виконання, персистентність, підвищення привілеїв, ухилення від захисту, доступ до облікових даних, розвідка, рух мережею, збір, C2, ексфільтрація, вплив) і техніками (способи досягнення). Для протидії ЗПК ATT\&CK слугує одночасно моделлю загроз (що може робити зловмисник), мовою опису поведінки зразків (розділ~\ref{chap:ch07}) і картою покриття захисту: зіставивши наявні засоби виявлення з матрицею технік, організація бачить «білі плями». Перевага ATT\&CK над абстрактними класифікаціями в тому, що вона побудована на \emph{реально спостережених} техніках і прив'язана до конкретних угруповань і зразків, що робить її практичним інструментом як для «червоних» команд (моделювання атак), так і для «синіх» (проєктування виявлення). Для кожної техніки база містить опис, приклади використання реальними групами, способи виявлення й пом'якшення, що дозволяє від моделі загроз одразу переходити до вибору контрзаходів (зв'язок із D3FEND~\cite{mitre-defend}). Важливе застереження: покриття технік не слід сприймати як просте «закривання клітинок матриці», бо різні техніки мають різну частоту, значущість і складність виявлення; пріоритизувати слід за реальною загрозою для конкретної організації, спираючись на розвідку про релевантні угруповання (розділ~\ref{chap:ch17}). Відображення типових технік ЗПК на тактики наведено в табл.~\thechapter.1.

\tabhead{Таблиця~\thechapter.1 --- Типові техніки злоякісного коду за тактиками ATT\&CK}

\begin{center}\footnotesize
\begin{tabularx}{\linewidth}{L{26mm} Y}
\toprule
Тактика & Приклади технік ЗПК \\
\midrule
Початковий доступ & фішинг із вкладенням, експлуатація публічного сервісу, ланцюг постачання \\
Виконання & запуск користувачем, скрипти, інтерпретатори, планувальник \\
Персистентність & автозапуск, служби, заплановані завдання, модифікація реєстру (розділ~\ref{chap:ch05}) \\
Ухилення від захисту & пакування, обфускація, вимкнення засобів, LotL (розділи~\ref{chap:ch08}, \ref{chap:ch13}) \\
Доступ до облікових даних & дамп пам'яті, кейлогінг, крадіжка токенів \\
Рух мережею & віддалене виконання, викрадені облікові дані \\
C2 & зашифрований канал, DNS, легітимні сервіси (розділ~\ref{chap:ch11}) \\
Вплив & шифрування (вимагач), знищення (вайпер), маніпуляція процесами \\
\bottomrule
\end{tabularx}
\end{center}

\section{Дерева атак}\label{sec:ch04-trees}

\term{Дерево атак}~\cite{ch04:schneier,ch04:kordy} --- ієрархічна модель, корінь якої --- мета зловмисника, листя --- елементарні дії, а внутрішні вузли типу «АБО» (достатньо однієї) та «І» (потрібні всі) описують альтернативи й послідовності. Атрибутам листя (ймовірність, вартість) приписують значення, що поширюються вгору: для «АБО» ймовірність $p_{\text{АБО}} = 1-\prod(1-p_k)$, для «І» --- $p_{\text{І}} = \prod p_k$. На рис.~\thechapter.2 наведено дерево для мети «виконання ЗПК на робочій станції» з трьома альтернативними шляхами доставки.

\begin{tikzfig}
\begin{tikzpicture}[font=\scriptsize,
  g/.style={draw,rounded corners=2pt,align=center,text width=20mm,minimum height=7mm,inner sep=1pt},
  l/.style={draw,align=center,text width=19mm,minimum height=8mm,inner sep=1pt,font=\tiny},
  >={Stealth[length=1.4mm]}]
\node[g] (G) at (4,0) {Виконання ЗПК (АБО)};
\node[g] (A1) at (0,-1.4) {Фішинг (І)};
\node[g] (A2) at (4,-1.4) {Вразливий сервіс};
\node[g] (A3) at (8,-1.4) {USB-носій};
\node[l] (a) at (-0.9,-2.9) {доставка листа 0,9};
\node[l] (b) at (1.1,-2.9) {запуск вкладення 0,3};
\node[l] (c) at (4,-2.9) {експлуатація 0,15};
\node[l] (d) at (8,-2.9) {підключення+запуск 0,1};
\foreach \p/\c in {G/A1,G/A2,G/A3,A1/a,A1/b,A2/c,A3/d} \draw (\p.south) -- (\c.north);
\end{tikzpicture}
\figcap{Рис.~\thechapter.2 --- Дерево атак для мети «виконання злоякісного коду на робочій станції»}
\end{tikzfig}

За наведеними ймовірностями $p_{A1} = 0{,}9\cdot0{,}3 = 0{,}27$, $p_{A2} = 0{,}15$, $p_{A3} = 0{,}1$, а для кореня $p_G = 1-(1-0{,}27)(1-0{,}15)(1-0{,}1) = 0{,}442$. Заходи захисту моделюють як зниження ймовірності відповідних листків: навчання користувачів і фільтрація вкладень знижують запуск вкладення з 0,3 до 0,05 ($p_{A1}=0{,}045$); своєчасне оновлення --- експлуатацію з 0,15 до 0,02; контроль USB --- з 0,1 до 0,01. Тоді $p_G = 1-(1-0{,}045)(1-0{,}02)(1-0{,}01) = 0{,}074$, тобто зниження майже в 6~разів. Дерево наочно показує, що найбільший внесок дає найімовірніший шлях (фішинг), тому пріоритет --- заходи проти нього.

Розширенням дерев атак є \term{дерева атак і захисту}~\cite{ch04:kordy}, у яких контрзаходи подано як окремі вузли, що взаємодіють з атакувальними: захід блокує атаку, атака може обходити захід, що моделюється вкладеними піддеревами. Такі моделі ближчі до реальності, бо відображають «гонку» між нападом і захистом: встановлення EDR змушує зловмисника застосувати техніки ухилення (розділ~\ref{chap:ch08}), що, своєю чергою, можуть бути виявлені поведінковим аналізом тощо. Перевага дерев --- наочність і можливість кількісного аналізу; обмеження --- вони погано відображають паралельність, час і стан системи, тому для складних сценаріїв застосовують графи атак, байєсові мережі чи імітаційне моделювання. Для протидії ЗПК дерева атак особливо корисні на етапі проєктування захисту, коли треба обґрунтувати, які рубежі й у якому порядку впроваджувати за обмеженого бюджету.

\section{Кількісне оцінювання ризику}\label{sec:ch04-quant}

\subsection{Ризик як сподівані втрати}

Ризик від ЗПК оцінюють як математичне сподівання річних втрат (ALE):

\begin{equation}
\mathrm{ALE} = \sum_{i} \mathrm{ARO}_i \cdot \mathrm{SLE}_i,\qquad \mathrm{SLE}_i = \mathrm{AV}_i \cdot \mathrm{EF}_i,
\label{eq:ch04-ale}
\end{equation}

де $\mathrm{ARO}_i$ --- річна частота реалізації $i$-го сценарію; $\mathrm{SLE}_i$ --- втрати від однієї події; $\mathrm{AV}_i$ --- вартість активу; $\mathrm{EF}_i$ --- фактор впливу. Методологія FAIR~\cite{ch04:fair} деталізує оцінки через частоту подій загрози, уразливість і первинні/вторинні втрати з поданням невизначеності розподілами й моделюванням Монте-Карло. Для ЗПК важливо, що $\mathrm{EF}$ різко залежить від типу: для шпигунського коду втрати --- вартість викрадених даних і репутації; для вимагача --- простій плюс можливий викуп; для вайпера --- $\mathrm{EF}\to1$ (повне знищення), що робить такі сценарії пріоритетними попри меншу частоту.

\subsection{Ймовірність компрометації ланцюга}

Оскільки атака є ланцюгом, імовірність успішної компрометації оцінюють добутком імовірностей проходження обов'язкових етапів:

\begin{equation}
P_{\text{комп}} = \prod_{j=1}^{m} p_j(X),\qquad p_j(X) = p_j^{0}\,(1-e_j),
\label{eq:ch04-chain}
\end{equation}

де $p_j^0$ --- імовірність успіху $j$-го етапу без заходів; $e_j\in[0;1]$ --- ефективність заходу на цьому етапі. Формула~\eqref{eq:ch04-chain} кількісно виражає силу ешелонованого захисту: навіть помірно ефективні заходи на кількох етапах у добутку дають низьку сумарну ймовірність. Наприклад, три етапи з $p_j^0 = 0{,}8$ і заходами з $e_j = 0{,}7$ дають $P_{\text{комп}} = (0{,}24)^3 \approx 0{,}014$ замість $0{,}512$ без заходів. Для самопоширюваних загроз оцінку доповнюють епідемічною моделлю (розділ~\ref{chap:ch03}): навіть за успішного первинного проникнення масштаб ураження визначається $R_0$, тому сегментація й швидка ізоляція (зменшення $R_0$ нижче~1) обмежують шкоду~\cite{ch04:staniford,ch04:zou}.

\subsection{Інтеграція з оцінкою вразливостей}

Імовірність експлуатації вразливостей, які ЗПК використовує для доставки, оцінюють поєднанням CVSS (технічна серйозність)~\cite{ch02:cvss} і EPSS (емпірична ймовірність експлуатації протягом 30~діб)~\cite{ch04:epss}. Це дозволяє пріоритизувати усунення вразливостей за реальним ризиком, а не лише за балом серйозності: вразливість із помірним CVSS, але високим EPSS (активно експлуатується) небезпечніша за «страшну» на папері, але не експлуатовану. Прогнозування майбутніх загроз на основі історичних даних і поведінкових шаблонів досліджують із застосуванням штучного інтелекту~\cite{ch04:ilienko}.

Невизначеність оцінок --- невід'ємна риса управління ризиками щодо ЗПК. Частоти атак, ймовірності етапів і розміри втрат відомі лише приблизно, тому точкові оцінки створюють хибне відчуття визначеності. Сучасна практика~\cite{ch04:fair} вимагає подавати величини як калібровані інтервали (наприклад, 90-відсоткові) і обчислювати розподіл сумарного ризику методом Монте-Карло, отримуючи не одне число, а криву перевищення втрат --- імовірність того, що річні втрати перевищать заданий рівень. Таку криву зіставляють із толерантністю до ризику, що її задає керівництво. Для протидії ЗПК це особливо важливо через «важкі хвости» розподілу втрат: рідкісні, але катастрофічні події (руйнівна атака, що паралізує організацію) домінують у математичному сподіванні, тому оцінювання лише за «типовим» сценарієм систематично занижує ризик. Український досвід (розділ~\ref{chap:ch01}) переконливо показав реальність таких хвостів.

\section{Пріоритизація заходів та економіка}\label{sec:ch04-economics}

Доцільність заходу оцінюють показником рентабельності інвестицій у безпеку:

\begin{equation}
\mathrm{ROSI} = \frac{\Delta\mathrm{ALE} - C}{C},
\label{eq:ch04-rosi}
\end{equation}

де $\Delta\mathrm{ALE}$ --- зниження сподіваних втрат завдяки заходу; $C$ --- його річна вартість. Захід виправданий за $\mathrm{ROSI}>0$. Оскільки $\Delta\mathrm{ALE}$ залежить від уже впроваджених заходів, їх оцінюють послідовно в порядку спадання ROSI. Модель Гордона--Лоба~\cite{ch04:gordonloeb} показує, що оптимальні інвестиції в захист не перевищують $\approx 37\,\%$ сподіваних втрат. Проте суто фінансові показники не охоплюють усього: обов'язкові за законодавством заходи (розділ~\ref{chap:ch02}) впроваджуються незалежно від ROSI; детективні заходи мають цінність, яку важко виразити через $\Delta\mathrm{ALE}$; для сценаріїв із загрозою безпеці людей застосовують принцип «настільки низько, наскільки практично досяжно».

Ключовим є \emph{розміщення} заходів уздовж ланцюга ураження~\eqref{eq:ch04-chain}: заходи на ранніх етапах (зменшення поверхні атаки, навчання, фільтрація) дешевші й запобігають більшості масових атак; заходи на пізніх етапах (виявлення C2, поведінковий аналіз, реагування) дорожчі, але необхідні проти цілеспрямованих загроз, що подолали ранні рубежі. Збалансований захист покриває весь ланцюг, а не концентрується на одному етапі.

Для практичного ухвалення рішень корисно розрізняти три типи заходів за їхньою функцією. \emph{Запобіжні} заходи знижують імовірність успіху етапу (ефективність $e_j$ у~\eqref{eq:ch04-chain}): оновлення, контроль застосунків, навчання. \emph{Детективні} заходи не запобігають, але виявляють дію й уможливлюють реагування: моніторинг, EDR, аналіз журналів. \emph{Коригувальні} заходи обмежують шкоду й відновлюють стан: сегментація, ізольовані резервні копії, план відновлення. Помилка концентрації лише на запобіганні є поширеною й небезпечною: за припущення компрометації (розділ~\ref{chap:ch01}) детективні й коригувальні заходи стають не менш важливими, бо визначають, наскільки швидко й з якими втратами організація виявить і подолає атаку, що неминуче подолає частину запобіжних рубежів. Кількісно це відображається в тому, що інвестиції лише в запобігання мають спадну віддачу (кожен наступний запобіжний рубіж дорожчий і дає менший приріст $e_j$), тоді як збалансований портфель запобіжних, детективних і коригувальних заходів мінімізує сумарний ризик за фіксованого бюджету. Саме тому сучасні рамки (NIST~CSF~2.0) явно виділяють функції виявлення, реагування й відновлення поряд із захистом (розділ~\ref{chap:ch02}).

\section{Кейс: модель загроз організації}\label{sec:ch04-case}

Побудуємо модель загроз для організації середнього розміру (фінансовий сектор) за методикою розділу.

\emph{Контекст і профіль порушника.} Релевантні: масові кіберзлочинці (вимагачі, банківські трояни), цілеспрямовані групи (крадіжка коштів), інсайдери. Активи: клієнтські дані, платіжні системи, доступність сервісів.

\emph{Ланцюги загроз.} S1 --- вимагач через фішинг (ARO~$=2$, SLE~$=500$~тис.); S2 --- банківський троян через заражений сайт (ARO~$=1$, SLE~$=800$~тис.); S3 --- APT через ланцюг постачання (ARO~$=0{,}2$, SLE~$=3000$~тис.); S4 --- деструктивна атака (ARO~$=0{,}1$, SLE~$=5000$~тис.). Сумарний $\mathrm{ALE}_0 = 1000+800+600+500 = 2900$~тис.

\emph{Заходи й пріоритизація.} C1 --- навчання + фільтрація пошти + контроль макросів ($C=150$; знижує S1, S2); C2 --- EDR з поведінковим виявленням ($C=300$; знижує всі); C3 --- ізольовані незмінні резервні копії + план відновлення ($C=200$; знижує EF для S1, S4); C4 --- контроль ланцюга постачання, SBOM ($C=120$; знижує S3). За~\eqref{eq:ch04-chain} і~\eqref{eq:ch04-ale} послідовний розрахунок дає найвищий ROSI для C1 (дешевий, б'є найчастіші вектори) і C3 (різко знижує EF катастрофічних сценаріїв), далі C2 і C4.

\emph{Реєстр ризиків} (фрагмент, табл.~\thechapter.2) фіксує сценарії, рівень і рішення. Для найруйнівніших сценаріїв (S4) рішення --- зниження через C3 незалежно від ROSI, оскільки йдеться про виживання організації. Кейс ілюструє, що модель загроз для ЗПК має поєднувати ймовірнісне оцінювання з безумовним захистом від катастрофічних сценаріїв, а пріоритет надавати заходам на ранніх етапах ланцюга й ізольованому відновленню. Питання атрибуції, важливе для реагування, розглядають у межах багаторівневих моделей~\cite{ch04:lungol}.

\tabhead{Таблиця~\thechapter.2 --- Фрагмент реєстру ризиків злоякісного коду}

\begin{center}\footnotesize
\begin{tabularx}{\linewidth}{Y c c L{30mm}}
\toprule
Сценарій & ALE, тис. & Рівень & Рішення \\
\midrule
Вимагач через фішинг & 1000 & високий & C1 + C2 + C3 \\
Банківський троян & 800 & високий & C1 + C2 \\
APT через ланцюг постачання & 600 & середній & C4 + C2 \\
Деструктивна атака (вайпер) & 500 & критичний & C3 (безумовно) + C2 \\
\bottomrule
\end{tabularx}
\end{center}

\section*{Висновки до розділу}
\addcontentsline{toc}{section}{Висновки до розділу}

Моделювання загроз і управління ризиками пов'язують нормативні вимоги й технічні засоби в систему обґрунтованих рішень. Для злоякісного коду загроза реалізується як багатоетапна операція, описувана ланцюгом ураження й базою ATT\&CK; ешелонований захист створює асиметрію на користь захисника, бо розрив будь-якого обов'язкового етапу зриває операцію~\eqref{eq:ch04-chain}. Дерева атак дозволяють кількісно аналізувати альтернативні шляхи доставки й ефект заходів, показуючи пріоритет протидії найімовірнішому вектору (фішингу). Кількісне оцінювання ризику через ALE~\eqref{eq:ch04-ale} має враховувати залежність фактора впливу від типу коду: деструктивні сценарії з $\mathrm{EF}\to1$ пріоритетні попри меншу частоту. Для самопоширюваних загроз оцінку доповнює епідемічна модель, а поєднання CVSS і EPSS уточнює ризик експлуатації вразливостей. Пріоритизація заходів через ROSI~\eqref{eq:ch04-rosi} з урахуванням розміщення вздовж ланцюга ураження, доповнена безумовним захистом від катастрофічних сценаріїв (ізольоване відновлення), становить основу раціональної програми протидії. Модель загроз має будуватися на релевантному профілі порушника, а не на абстрактному ландшафті.

\selfcheck
\begin{enumerate}
\item Опишіть процес управління ризиками щодо ЗПК. Чому визначення релевантного профілю порушника є ключовим на етапі контексту?
\item Опишіть етапи класичного ланцюга ураження (рис.~\thechapter.1). Чому розрив будь-якого етапу зриває операцію і яка асиметрія з цього випливає?
\item Чим уніфікований ланцюг ураження доповнює класичний? Опишіть цикл внутрішньомережевої фази.
\item Як база ATT\&CK слугує одночасно моделлю загроз, мовою опису поведінки й картою покриття захисту (табл.~\thechapter.1)?
\item Побудуйте дерево атак для мети доставки ЗПК і обчисліть ймовірність досягнення кореня. Як моделюють ефект заходів?
\item За деревом рис.~\thechapter.2 перерахуйте $p_G$, якщо навчання знижує запуск вкладення до 0,1, а решта незмінна.
\item Запишіть формулу ALE~\eqref{eq:ch04-ale}. Чому фактор впливу різко залежить від типу ЗПК і чому деструктивні сценарії пріоритетні?
\item Поясніть формулу ймовірності компрометації ланцюга~\eqref{eq:ch04-chain}. Обчисліть $P_{\text{комп}}$ для чотирьох етапів із $p_j^0=0{,}7$ і $e_j=0{,}6$.
\item Як поєднання CVSS і EPSS уточнює пріоритизацію усунення вразливостей? Наведіть приклад.
\item Поясніть показник ROSI~\eqref{eq:ch04-rosi} і результат Гордона--Лоба. Які заходи впроваджують незалежно від ROSI?
\item Чому важливе розміщення заходів уздовж ланцюга ураження? Порівняйте заходи на ранніх і пізніх етапах.
\item Чому для найруйнівніших сценаріїв рішення ухвалюють незалежно від ROSI? Проілюструйте на кейсі підрозділу~\ref{sec:ch04-case}.
\item Як епідемічна модель (розділ~\ref{chap:ch03}) доповнює оцінку ризику для самопоширюваних загроз?
\item Дослідницьке: побудуйте повну модель загроз для організації обраного сектору з деревами атак, кількісним оцінюванням і реєстром ризиків, та обґрунтуйте пріоритизацію заходів.
\item Дослідницьке: розробіть методику зіставлення наявних засобів виявлення з матрицею ATT\&CK для оцінювання покриття захисту й виявлення «білих плям», застосуйте її до модельної інфраструктури.
\end{enumerate}
